\documentclass[10pt,aspectratio=169]{beamer}

% All the boilerplate is in deslides.sty
\usepackage{deslides}

\author{Ji\v{r}\'i Lebl}

\institute[OSU]{%
Oklahoma State University%
%Departemento pri Matematiko de Oklahoma {\^S}tata Universitato%
}

\title{29. Power series, part 1\\(Notes on Diffy Qs, 7.1)}

\date{}

\begin{document}

\begin{frame}
\titlepage

%\bigskip

\begin{center}
The textbook: \url{https://www.jirka.org/diffyqs/}

or \url{https://jirilebl.github.io/diffyqs/}
\end{center}
\end{frame}

\begin{frame}
Wouldn't it be nice if all functions were polynomials?

E.g., we could just apply undetermined coefficients all the time.

\pause
\medskip

Instead of polynomials, consider \emph{power series}:
\[
\sum_{k=0}^\infty a_k {(x-x_0)}^k =
a_0 + 
a_1 (x-x_0) +
a_2 {(x-x_0)}^2 +
a_3 {(x-x_0)}^3 + \cdots.
\]
$a_0,a_1,a_2,\ldots,a_k,\ldots$ (coefficients) and $x_0$ (center) are
constants.

\pause
\medskip

The \emph{partial sum} is the sum of the first $n$ terms:
\[
S_n(x)
=
\sum_{k=0}^n a_k {(x-x_0)}^k
=
a_0 + a_1 (x-x_0) + a_2 {(x-x_0)}^2 + a_3 {(x-x_0)}^3 + \cdots + a_n {(x-x_0)}^n ,
\]

\pause
Say the series \emph{converges at $x$} if the following limit exists:
\[
\sum_{k=0}^\infty a_k {(x-x_0)}^k = 
\lim_{n\to \infty} S_n(x) = \lim_{n\to\infty} \sum_{k=0}^n a_k {(x-x_0)}^k
\]
\end{frame}

\begin{frame}
$\displaystyle \sum_{k=0}^\infty a_k {(x-x_0)}^k$
\quad
always converges at $x=x_0$ (to $a_0$).

\pause
\medskip

The series is \emph{convergent}
if it converges at some $x \neq x_0$ (\emph{divergent} otherwise).

\pause
\medskip

\textbf{Example:}
\quad
$\displaystyle
\sum_{k=0}^\infty \frac{1}{k!} x^k = 
1 + x + \frac{x^2}{2} + \frac{x^3}{6} + \cdots$
\quad
converges at every $x$ (to $e^x$ in fact).

\pause
\vspace*{-8pt}

\hfill
{\small (Recall: $k! = 1\cdot 2\cdot 3 \cdots k$ and $0! = 1$.)}


\pause
The series 
\emph{converges absolutely} at $x$ if
\quad
$\displaystyle
\sum_{k=0}^\infty
\lvert a_k \rvert \, {\lvert x-x_0 \rvert}^k 
$
\quad
converges $\Bigl($ {\small $\!\displaystyle \lim\limits_{n\to\infty}
\sum\limits_{k=0}^n
\lvert a_k \rvert \, {\lvert x-x_0 \rvert}^k$ exists}
{$\!\Bigr)$}

\pause
absolute convergence $\Rightarrow$ convergence, but
not vice-versa.

\pause
\textbf{Example:}
\quad
$\displaystyle
\sum_{k=1}^\infty \frac{1}{k} x^k
$\quad
converges absolutely whenever $-1 < x < 1$.

\vspace*{-7pt}

\pause
It converges at $x=-1$:
By the alternating series test,
$\displaystyle \sum_{k=1}^\infty \frac{{(-1)}^k}{k}$ converges,

\vspace*{-7pt}

\pause
but not absolutely as $\displaystyle \sum_{k=1}^\infty \frac{1}{k}$ does not converge.

\vspace*{-7pt}

\pause
\hfill That also means that the series diverges at $x=1$.

\end{frame}

\begin{frame}

Suppose a power series converges absolutely at some $x_1$
and $x$ is no further than $x_1$ to $x_0$:

\medskip

\qquad
$\lvert x - x_0  \rvert \leq \lvert x_1 - x_0 \rvert$
\pause
\qquad
$\Rightarrow$
\qquad
$\lvert a_k \rvert \, {\lvert x-x_0\rvert }^k \leq \lvert a_k \rvert
\, {\lvert x_1-x_0 \rvert}^k$ \quad for all $k$.

\pause
\medskip

\qquad
$\phantom{\lvert x - x_0  \rvert \leq \lvert x_1 - x_0 \rvert}$
\qquad
$\Rightarrow$
\qquad
$\displaystyle
\sum_{k=0}^\infty \lvert a_k \rvert \, {\lvert x-x_0\rvert }^k$
\quad
converges ~~{\small (a sum of smaller numbers)}

\pause
\medskip

\textbf{Theorem:}
For a power series, there exists a number
$\rho$ (we allow $\rho=\infty$)
called the \emph{radius of convergence} such that
the series converges absolutely on the interval
$(x_0-\rho,x_0+\rho)$ and diverges for $x < x_0-\rho$ and $x > x_0+\rho$.
We write $\rho=\infty$ if
the series converges for all $x$.

\begin{center}
\scalebox{0.7}{\subimport*{../figures/}{ps-conv.pdf_t}}
\end{center}

\pause
\vspace*{-5pt}

For
\quad
$\displaystyle
\sum_{k=0}^\infty \frac{1}{k!} x^k$
\quad
$\rho = \infty$.
\qquad
\pause
For
\quad
$\displaystyle
\sum_{k=0}^\infty \frac{1}{k} x^k$
\quad
$\rho = 1$.

\pause
\vspace*{-5pt}

$\rho = 0$ if the series is
divergent.
\qquad
E.g.,
\quad
$\displaystyle
\sum_{k=0}^\infty k! x^k$
\quad
diverges for all $x \neq 0$,
so $\rho=0$.

\end{frame}

\begin{frame}
How do we find $\rho$?  \pause
We use the \emph{ratio test} or the \emph{root test}.

\pause
\textbf{Ratio test:}
Given
\quad
$\displaystyle
\sum_{k=0}^\infty c_k
$,
\quad
suppose
\quad
$L = \lim_{k\to\infty} \left \lvert \frac{c_{k+1}}{c_k} \right \rvert$.

\pause
$\Rightarrow$ \quad
the series converges absolutely if $L < 1$ and
diverges if $L > 1$.

\pause
\medskip

Given
\quad
$\displaystyle
\sum_{k=0}^\infty a_k {(x-x_0)}^k
$,
\quad
$
\displaystyle
L = \lim_{k\to\infty} \left \lvert \frac{c_{k+1}}{c_k} \right \rvert
\pause
=
\lim_{k\to\infty} \left \lvert
\frac{a_{k+1} {(x - x_0)}^{k+1}}{a_k {(x - x_0)}^k}
\right \rvert
=
\pause
\lim_{k\to\infty} \left \lvert
\frac{a_{k+1}}{a_k}
\right \rvert
\lvert  x - x_0 \rvert
$.

Let
\quad
$\displaystyle A =
\lim_{k\to\infty} \left \lvert
\frac{a_{k+1}}{a_k}
\right \rvert$.
\quad
\pause
The series converges absolutely if $1 > L = A \lvert x - x_0 \rvert$.

\pause
\medskip

The radius of convergence is $\displaystyle \rho = \nicefrac{1}{A}$
\qquad
($A = 0 \, \Rightarrow \, \rho=\infty$ and
$A = \infty \, \Rightarrow \, \rho=0$).

\pause
\medskip

\textbf{Root test:} Same but
\quad
$\displaystyle
L = \lim_{k\to\infty} \sqrt[k]{\lvert c_k \rvert}$
\quad
and so
\quad
$\displaystyle
A = \lim_{k\to\infty} \sqrt[k]{\lvert a_k \rvert}$.

\end{frame}

\begin{frame}


\textbf{Example:}
Consider
\quad
$\displaystyle
\sum_{k=0}^\infty 2^{-k} {(x-1)}^k$.

\pause
Compute
\quad
$\displaystyle
A = \lim_{k\to\infty} 
\left \lvert
\frac{a_{k+1}}{a_k}
\right \rvert
\pause
=
\lim_{k\to\infty} 
\left \lvert
\frac{2^{-k-1}}{2^{-k}}
\right \rvert
\pause
=
\lim_{k\to\infty} 
2^{-1} = \frac{1}{2}$.

\pause
\medskip

$\Rightarrow$ $\rho = 2$.
\pause
As $x_0 = 1$,
the series converges absolutely on
the interval $(1-2,1+2) = (-1,3)$.

\medskip
\pause

Root test would also work:
\quad
$\displaystyle
A = 
\lim_{k\to\infty} 
\sqrt[k]{\lvert
a_k
\rvert}
=
\lim_{k\to\infty} 
\sqrt[k]{\lvert
2^{-k}
\rvert}
=
\lim_{k\to\infty} 
2^{-1}
=
\frac{1}{2}$.

\pause
\medskip

\textbf{Example:}
Consider
\quad
$\displaystyle
\sum_{k=1}^\infty \frac{1}{k^k} {x}^k$.

\pause
Compute
\quad
$\displaystyle
A =
\lim_{k\to\infty} 
\sqrt[k]{\lvert a_k \rvert}
\pause
=
\lim_{k\to\infty} 
\sqrt[k]{
\left\lvert\frac{1}{k^k}\right\rvert}
=
\pause
\lim_{k\to\infty} 
\sqrt[k]{
{\left\lvert\frac{1}{k}\right\rvert}^{k}}
\pause
=
\lim_{k\to\infty} 
\frac{1}{k}
=
0$.

\pause
\medskip

The radius of convergence is $\infty$: The series
converges everywhere.

\end{frame}

\begin{frame}
The limits
$\bigl \lvert \frac{a_{k+1}}{a_k} \bigr \rvert$
or
$\sqrt[k]{\lvert a_k \rvert}$
might not exist always.

\pause
\medskip

Sometimes we can't compute the $A$, but we can compute the
$L$ for the terms given.

\pause
\medskip

\textbf{Example:}
Consider \quad
$\displaystyle \sum_{k=1}^\infty 2^k x^{2k}$.

\pause
\medskip

Note that $a_n = 0$ if $n$ is odd and $a_n = a_{2k} = 2^k = 2^{n/2}$
if $n$ is even.

\medskip
\pause

We cannot compute $A$:
$\bigl \lvert \frac{a_{n+1}}{a_n} \bigr \rvert$
is undefined if $n$ is odd and
$\sqrt[n]{\lvert a_n \rvert}$ jumps between $0$
and $\sqrt{2}$.

\pause
\medskip

But
\quad
$\displaystyle
L = \lim_{k\to \infty}
\sqrt[k]{\lvert 2^k x^{2k} \rvert}
= 2 {\lvert x\rvert}^2$

\pause
\medskip

The series converges absolutely if \quad
$L = 2 {\lvert x\rvert}^2 < 1$,

\pause
\medskip

in other words, if
\quad ${\lvert x\rvert}^2 < \dfrac{1}{2}$
\quad
or
\quad ${\lvert x\rvert} < \dfrac{1}{\sqrt{2}}$.

\pause
\medskip

The radius of convergence is $\rho = \dfrac{1}{\sqrt{2}}$.

\end{frame}

\begin{frame}

At the endpoints,
$x = x_0-\rho$ and $x = x_0+\rho$,
a power series may or may not converge.

\pause
\medskip

Neither ratio nor root test are going to apply at the
endpoints.

\pause
\medskip

Sometimes convergence at the endpoints is important.

\pause
\medskip

Fortunately, for our purposes, the radius of convergence is enough.

\end{frame}


\end{document}
